\documentclass[12pt]{article}

\usepackage{booktabs}
\usepackage{tabularx}
\usepackage{threeparttable}
\usepackage{makecell}
\usepackage{amsmath}
\usepackage{amssymb}
\usepackage{geometry}
\usepackage{setspace}
\usepackage{parskip}

\geometry{margin=1in}
\onehalfspacing

\title{Preliminary Results\\
\large Causal Effect of Social Media Sentiment on Stock Price Movements}
\date{\today}

\begin{document}

\maketitle

\section{Analytical Approach}

This study seeks to estimate the causal effect of social media sentiment on stock prices using two complementary identification strategies.

The first employs a firm-level fixed effects regression using \texttt{pyfixest}. This allows for the control of time-invariant factors within each firm, to account for shocks to individual firms which may skew results. This strategy assumes that confounders influence the outcome in some linear manner, making it the simpler of the two models.

The second strategy employs a doubly debiased machine learning estimator --- DoubleML Partial Linear Regression (\texttt{DoubleMLPLR}) combined with \texttt{XGBoost}, a nuisance learner. This strategy allows XGBoost to identify connections between the treatment and dependent variables across the matrix of confounders to generate predictive models, then using DoubleMLPLR to identify a coefficient, producing a treatment effect. For this initial set of estimations, I opted for 5-fold cross-fitting done over 20 repetitions, with 1000 bootstrap draws to arrive at an estimator.

With both methods, data consisted of daily panel data pulled from Bloomberg. Sentiment ratings were provided using Bloomberg's proprietary tool as well.

\newpage
\section{Fixed Effects Results}

Table~\ref{tab:fe_main} presents the four FE-OLS specifications. None of the models produced statistically significant results for Twitter sentiment's impact on stock price. This indicates that interactions between tweets and traders do not influence behavior more than other factors. However, each model provides interesting insights into the relationship between social media and share price.

Model 1 is the baseline regression of return on \texttt{twitter\_sent} without confounders. The estimated coefficient is $-$0.370, indicating that a one-unit increase in Twitter sentiment is associated with a 0.37 cent decrease in the same-day return. While the sign is consistent with expected findings, the result is not statistically significant. Also, the $R^2$ (0.002) tells us that sentiment explains little of what is driving share prices.

Model 2 adds the full confounder matrix: intraday price range (\texttt{px\_high}, \texttt{px\_low}), market capitalization, total equity, debt-to-equity ratio, trading volume, news sentiment, RSI (30-day), and the 50-day moving average. The \texttt{twitter\_sent} coefficient's significance improved slightly; however, it is still insignificant. However, \texttt{news\_sent} is substantial and significant. RSI was significant as well. This tells us that news sentiment and RSI are stronger predictors of returns than is Twitter sentiment, at least in an OLS framework.

Models 3 and 4 were intended to check other approaches to see if they yielded better results. Model 3 uses a negative sentiment measure \texttt{neg\_twitter\_sent}, which is zero on positive-sentiment days and equals the raw score on negative-sentiment days. The coefficient was $-$0.072, smaller in magnitude than Model 1 and not statistically significant. This suggests that negative sentiment does not have particularly more effect than the entire spectrum of sentiment scores.

Model 4 borrowed from a previous study and used tweet counts (\texttt{twitter\_neg\_count}) as the treatment. The coefficient is 0.000533, indicating that to some extent any press is good press, but not in a statistically significant way.

\begin{table}[htbp]
\centering
\caption{Fixed Effects OLS --- Twitter Sentiment and Stock Returns}
\label{tab:fe_main}
\begin{threeparttable}
\begingroup
\renewcommand\cellalign{t}
\renewcommand\arraystretch{1}
\setlength{\tabcolsep}{3pt}
\begin{tabularx}{\linewidth}{@{}>{\raggedright\arraybackslash}l>{\centering\arraybackslash}X>{\centering\arraybackslash}X>{\centering\arraybackslash}X>{\centering\arraybackslash}X}
\toprule
 & \multicolumn{4}{c}{return} \\
\cmidrule(lr){2-5}
 & (1) & (2) & (3) & (4) \\
\midrule
\addlinespace[1ex]
twitter\_sent & \makecell{-0.37 \\ (0.23)} & \makecell{-0.422 \\ (0.325)} &  &  \\
\addlinespace[0.5ex]
px\_high &  & \makecell{-0.085 \\ (0.121)} &  &  \\
\addlinespace[0.5ex]
px\_low &  & \makecell{0.091 \\ (0.122)} &  &  \\
\addlinespace[0.5ex]
mkt\_cap &  & \makecell{0 \\ (0.000001)} &  &  \\
\addlinespace[0.5ex]
total\_equity &  & \makecell{0.000001 \\ (0.000004)} &  &  \\
\addlinespace[0.5ex]
debt\_to\_equity &  & \makecell{-0.000089 \\ (0.000049)} &  &  \\
\addlinespace[0.5ex]
volume &  & \makecell{0 \\ (0)} &  &  \\
\addlinespace[0.5ex]
news\_sent &  & \makecell{-1.056 \\ (0.114)} &  &  \\
\addlinespace[0.5ex]
rsi\_30 &  & \makecell{0.084 \\ (0.016)} &  &  \\
\addlinespace[0.5ex]
ma\_50 &  & \makecell{-0.013 \\ (0.011)} &  &  \\
\addlinespace[0.5ex]
twitter\_neg\_count &  & \makecell{0.00085 \\ (0.001)} &  & \makecell{0.000533 \\ (0.000971)} \\
\addlinespace[0.5ex]
neg\_twitter\_sent &  &  & \makecell{-0.072 \\ (0.271)} &  \\
\addlinespace[0.5ex]
\midrule
\addlinespace[1ex]
ticker FE & $\times$ & $\times$ & $\times$ & $\times$ \\
\midrule
\addlinespace[1ex]
Observations & 160,103 & 147,128 & 160,103 & 159,232 \\
\addlinespace[0.5ex]
$R^2$ & 0.002 & 0.019 & 0.002 & 0.002 \\
\addlinespace[0.5ex]
\bottomrule
\end{tabularx}
\endgroup
\begin{tablenotes}
\footnotesize
\item Heteroskedasticity-robust SE in parentheses. Company fixed effects in all models. Model 3 treatment is \texttt{clip(twitter\_sent, max=0)}. Model 4 per Teti et al.\ (2019).
\end{tablenotes}
\end{threeparttable}
\end{table}

\newpage
\section{DoubleML Results --- Sentiment}

Table~\ref{tab:dml_twitter} presents the DoubleMLPLR estimates, and the contrast with the FE results is stark.

Run 1, using \texttt{twitter\_sent} as the continuous treatment, yielded a coefficient of $-$0.924 with strong significance. This indicates that changes in Twitter sentiment negatively influence share prices, regardless of the direction of sentiment. So perhaps all Twitter coverage is bad press. However, in Run 2, which uses \texttt{twitter\_neg\_count} as the treatment to determine if tweet volume played a role, the DoubleML estimate is $+$0.003987. This positive, significant coefficient indicates that a higher count of negative tweets is associated with higher same-day returns. This was unexpected, and means that potentially tweet volume and tweet sentiment are not necessarily connected. My interpretation is that when a company is focused upon, the total number of tweets for the day increases, also increasing the amount of negative comments, but not impacting sentiment across the day.

\begin{table}[htbp]
\centering
\caption{DoubleML PLR --- Twitter Sentiment and Stock Returns}
\label{tab:dml_twitter}
\begin{tabular}{lcc}
\hline\hline
 & (1) Twitter Sentiment & (2) Neg Tweet Count \\
\hline
twitter\_sent & $-$0.924313*** & --- \\
& (0.225639) &  \\
twitter\_neg\_count & --- & 0.003987*** \\
&  & (0.001503) \\
\hline
Estimator & DoubleML PLR & DoubleML PLR \\
ML Learner & XGBoost (GPU) & XGBoost (GPU) \\
\hline\hline
\multicolumn{3}{p{0.9\linewidth}}{\footnotesize \textit{Notes:} Heteroskedasticity-robust SE in parentheses. XGBoost learner with 5-fold cross-fitting, 20 repetitions, 1000 bootstrap draws. Run 2 per Teti et al.\ (2019): polarity-broken tweet counts outperform total count. *** $p<0.01$, ** $p<0.05$, * $p<0.1$.}
\end{tabular}
\end{table}

Table~\ref{tab:dml_news} shows results from applying the DoubleML framework to news sentiment. Run 1 shows that when using \texttt{news\_sent} as treatment, we get a confident estimator of $-$0.759. This result is consistent with FE Model 2 (which showed news sentiment as the strongest predictor of return), but adds additional credibility. My second run for news sentiment was incorrectly set up and is being removed, as it is essentially the same as the run with negative tweet counts. This is from a limitation in Bloomberg's data, wherein they cannot collect a count of news mentions across a day.

\begin{table}[htbp]
\centering
\caption{DoubleML PLR --- News Sentiment and Stock Returns}
\label{tab:dml_news}
\begin{tabular}{lcc}
\hline\hline
 & (1) News Sentiment & (2) Neg News Volume \\
\hline
news\_sent & $-$0.759391*** & --- \\
& (0.134525) &  \\
neg\_news\_vol & --- & 0.000588** \\
&  & (0.000232) \\
\hline
Estimator & DoubleML PLR & DoubleML PLR \\
ML Learner & XGBoost (GPU) & XGBoost (GPU) \\
\hline\hline
\multicolumn{3}{p{0.9\linewidth}}{\footnotesize \textit{Notes:} Heteroskedasticity-robust SE in parentheses. XGBoost learner with 5-fold cross-fitting, 20 repetitions, 1000 bootstrap draws. Run 2 treatment: $\texttt{neg\_news\_vol} = \texttt{twitter\_count} \times \mathbf{1}(\texttt{news\_sent} < 0)$. *** $p<0.01$, ** $p<0.05$, * $p<0.1$.}
\end{tabular}
\end{table}

\newpage
\section{DoubleML Results --- Intraday Price High}

Table~\ref{tab:dml_px_high} shows results for the examination of \texttt{px\_high} (highest sale price of the day) as the outcome, while leaving the confounder set the same as in other runs. Run 1, \texttt{twitter\_sent} $\sim$ \texttt{px\_high}, yielded $+$1.610. This positive coefficient means positive Twitter sentiment raises the intraday price ceiling --- a result in direct contrast with the negative \texttt{twitter\_sent} $\sim$ \texttt{return} estimate. Then in Run 2, \texttt{twitter\_neg\_count} $\sim$ \texttt{px\_high} produced a significant estimator of $-$0.002082. Again, this is the opposite direction from the previous runs. This indicates that either \texttt{px\_high} is being influenced differently than \texttt{return}, or that factors impacting highs resolve themselves over the course of the trading day.

\begin{table}[htbp]
\centering
\caption{DoubleML PLR --- Sentiment and Intraday Price High (\texttt{px\_high})}
\label{tab:dml_px_high}
\begin{tabular}{lcc}
\hline\hline
 & (1) Twitter Sentiment & (2) Neg Tweet Count \\
\hline
twitter\_sent & $+$1.610130*** & --- \\
& (0.217458) &  \\
twitter\_neg\_count & --- & $-$0.002082* \\
&  & (0.001189) \\
\hline
Estimator & DoubleML PLR & DoubleML PLR \\
ML Learner & XGBoost (GPU) & XGBoost (GPU) \\
\hline\hline
\multicolumn{3}{p{0.9\linewidth}}{\footnotesize \textit{Notes:} Heteroskedasticity-robust SE in parentheses. XGBoost learner with 5-fold cross-fitting, 20 repetitions, 1000 bootstrap draws. Outcome is intraday high price (\texttt{px\_high}); \texttt{return} excluded from confounders. Run 2 per Teti et al.\ (2019). *** $p<0.01$, ** $p<0.05$, * $p<0.1$.}
\end{tabular}
\end{table}

\newpage
\section{Robustness Checks}

\textbf{Placebo test.} To test whether today's Twitter sentiment predicts yesterday's return, I regress \texttt{lag1} (prior-day return) on \texttt{twitter\_sent} within firms. If there were no reverse causality, the coefficient should be statistically indistinguishable from zero. However, the estimated coefficient is $+$2.348, making it both large and highly significant. This is a serious concern because it indicates that Twitter sentiment is more likely caused by returns than the other way around.

\textbf{No-reversal test.} I estimate all sentiment lags jointly (lags 1, 2, 3, 5, and 7) to test whether only lag 1 is significant. This would be expected if sentiment reflects news that is fully absorbed on day 1. None of the lag coefficients are statistically significant, and the signs alternate. This result is inconclusive, and the lack of significance at lag 1 suggests that any causal sentiment signal is likely too weak to survive multicollinear estimation against lags 2--7.

\begin{table}[htbp]
\centering
\caption{FE Robustness --- Placebo and No-Reversal Tests}
\label{tab:fe_robustness}
\begin{threeparttable}
\begingroup
\renewcommand\cellalign{t}
\renewcommand\arraystretch{1}
\setlength{\tabcolsep}{3pt}
\begin{tabularx}{\linewidth}{@{}>{\raggedright\arraybackslash}l>{\centering\arraybackslash}X>{\centering\arraybackslash}X}
\toprule
 & \multicolumn{1}{c}{lag1} & \multicolumn{1}{c}{return} \\
\cmidrule(lr){2-2} \cmidrule(lr){3-3}
 & (1) & (2) \\
\midrule
\addlinespace[1ex]
twitter\_sent & \makecell{2.348 \\ (0.225)} &  \\
\addlinespace[0.5ex]
sent\_lag1 &  & \makecell{$-$0.322 \\ (0.196)} \\
\addlinespace[0.5ex]
sent\_lag2 &  & \makecell{0.088 \\ (0.196)} \\
\addlinespace[0.5ex]
sent\_lag3 &  & \makecell{$-$0.225 \\ (0.192)} \\
\addlinespace[0.5ex]
sent\_lag5 &  & \makecell{$-$0.285 \\ (0.175)} \\
\addlinespace[0.5ex]
sent\_lag7 &  & \makecell{0.222 \\ (0.166)} \\
\addlinespace[0.5ex]
\midrule
\addlinespace[1ex]
ticker FE & $\times$ & $\times$ \\
\midrule
\addlinespace[1ex]
Observations & 159,606 & 156,589 \\
\addlinespace[0.5ex]
$R^2$ & 0.003 & 0.002 \\
\addlinespace[0.5ex]
\bottomrule
\end{tabularx}
\endgroup
\begin{tablenotes}
\footnotesize
\item Placebo: \texttt{lag1} $\sim$ \texttt{twitter\_sent} $|$ ticker. No-reversal: all sentiment lags estimated jointly. Only \texttt{sent\_lag1} should be significant under a causal interpretation (Gu \& Kurov, 2020).
\end{tablenotes}
\end{threeparttable}
\end{table}

\textbf{Lag decay.} Table~\ref{tab:lag_decay} estimates each lag separately in a univariate DoubleML specification. The results fail the no-reversal test comprehensively: all lagged sentiment variables --- at lags 1, 2, 3, 5, and 7 --- are highly significant ($p < 0.001$) and negative for both Twitter and news sentiment. Because there is no decay pattern, this means that rises in price are likely due to factors other than tweets or news (being the result of more permanent information or performance), or that reverse causality is in effect and price changes are driving social media sentiment.

\begin{table}[htbp]
\centering
\caption{Lag Decay --- Lagged Sentiment and Current Return (DoubleML PLR)}
\label{tab:lag_decay}
\begin{tabular}{llcccc}
\hline\hline
Treatment & Lag & Coefficient & Std Err & $p$-value & 95\% CI \\
\hline
news\_sent    & 1 & $-$1.169012*** & & 0.0000 & [$-$1.446810, $-$0.888520] \\
twitter\_sent & 1 & $-$1.215320*** & & 0.0000 & [$-$1.520064, $-$0.910576] \\
news\_sent    & 2 & $-$0.793453*** & & 0.0000 & [$-$1.037684, $-$0.543157] \\
twitter\_sent & 2 & $-$0.957274*** & & 0.0000 & [$-$1.243073, $-$0.672875] \\
news\_sent    & 3 & $-$1.363310*** & & 0.0000 & [$-$1.551271, $-$1.177750] \\
twitter\_sent & 3 & $-$1.122099*** & & 0.0000 & [$-$1.411447, $-$0.832751] \\
news\_sent    & 5 & $-$0.942506*** & & 0.0000 & [$-$1.129350, $-$0.755663] \\
twitter\_sent & 5 & $-$1.126149*** & & 0.0000 & [$-$1.396626, $-$0.856494] \\
news\_sent    & 7 & $-$0.638419*** & & 0.0000 & [$-$0.817363, $-$0.459343] \\
twitter\_sent & 7 & $-$0.674264*** & & 0.0000 & [$-$0.960148, $-$0.385289] \\
\hline
Estimator & \multicolumn{5}{l}{DoubleML PLR, XGBoost (GPU), 5-fold, 20 reps} \\
\hline\hline
\multicolumn{6}{p{0.95\linewidth}}{\footnotesize \textit{Notes:} Each lag estimated separately (univariate). Outcome: current-day return. Confounder set excludes contemporaneous sentiment to avoid endogeneity. XGBoost learner, 5-fold, 20 reps, 1000 bootstrap draws. *** $p<0.01$, ** $p<0.05$, * $p<0.1$.}
\end{tabular}
\end{table}

\textbf{Impact persistence.} Table~\ref{tab:impact_persist} estimates the effect of current sentiment on future returns. Under the hypothesis that tweet effects decay quickly (within 2--3 days) while news effects persist, \texttt{twitter\_sent} should be most significant at lead 1 and insignificant thereafter. The results show it insignificant at every lead. While \texttt{news\_sent} produces significant estimators at several leads, the signs are inconsistent. Taken together, this tells us that while Twitter sentiment is likely moot in the eyes of traders, news varies based on how the specific content matters to traders --- likely, important information is conveyed that does not tip the sentiment scales highly but is of value for determining tradability of stocks (e.g., exposure to a new regulation in a given industry).

\begin{table}[htbp]
\centering
\caption{Impact Persistence --- Current Sentiment and Future Return (DoubleML PLR)}
\label{tab:impact_persist}
\begin{tabular}{llcccc}
\hline\hline
Treatment & Lead & Coefficient & Std Err & $p$-value & 95\% CI \\
\hline
news\_sent    & 1 & $-$0.357017** &  & 0.0186 & [$-$0.657748, $-$0.059756] \\
twitter\_sent & 1 & $-$0.111458   &  & 0.6094 & [$-$0.537107,  $+$0.316047] \\
news\_sent    & 2 & $+$0.311669*** &  & 0.0090 & [$+$0.074430,  $+$0.545631] \\
twitter\_sent & 2 & $+$0.217330   &  & 0.3008 & [$-$0.194791,  $+$0.628997] \\
news\_sent    & 3 & $-$0.226585*  &  & 0.0575 & [$-$0.457798,  $+$0.007184] \\
twitter\_sent & 3 & $-$0.071248   &  & 0.7282 & [$-$0.480987,  $+$0.330521] \\
news\_sent    & 5 & $-$0.215964*  &  & 0.0578 & [$-$0.438797,  $+$0.007146] \\
twitter\_sent & 5 & $-$0.033204   &  & 0.8668 & [$-$0.420425,  $+$0.354753] \\
news\_sent    & 7 & $+$0.342398*** &  & 0.0026 & [$+$0.112941,  $+$0.565643] \\
twitter\_sent & 7 & $+$0.027558   &  & 0.8938 & [$-$0.375978,  $+$0.432329] \\
\hline
Estimator & \multicolumn{5}{l}{DoubleML PLR, XGBoost (GPU), 5-fold, 20 reps} \\
\hline\hline
\multicolumn{6}{p{0.95\linewidth}}{\footnotesize \textit{Notes:} Each lead estimated separately. Outcome: \texttt{return\_lead}$n$ (return $n$ days ahead). Hypothesis: \texttt{twitter\_sent} decays by lead 2--3; \texttt{news\_sent} persists to lead 5--7 (Gu \& Kurov, 2020). XGBoost learner, 5-fold, 20 reps, 1000 bootstrap draws. *** $p<0.01$, ** $p<0.05$, * $p<0.1$.}
\end{tabular}
\end{table}

\newpage
\section{Teti et al.\ (2019) Replication}

Teti et al.\ (2019) examine whether Twitter sentiment predicts open-to-close stock returns using a panel of S\&P 500 firms, with a key focus on heterogeneous effects by firm size. Their central finding is that large-cap firms --- which attract disproportionate social media attention --- exhibit stronger sentiment--return linkages than small-cap firms. To replicate this, I run three sets of models using Bloomberg's daily Twitter sentiment and count data.

\textbf{SET 1 --- Pooled OLS.} Table~\ref{tab:teti_set1} presents the pooled OLS specifications. Model 1 (baseline Twitter sentiment only) produces a small negative coefficient of $-$0.013, which is not statistically significant. This aligns with the FE results in Section 2. Model 3, which adds news sentiment as a co-treatment, shows \texttt{news\_sent} is strongly significant at $-$0.122 ($p < 0.01$), while \texttt{twitter\_sent} remains flat. Model 4 includes both simultaneously: the pattern holds. News sentiment dominates; Twitter sentiment adds nothing beyond it. This is a consistent theme across all specifications so far.

\begin{table}[htbp]
\centering
\caption{Teti et al.\ (2019) Replication --- SET 1: Pooled OLS}
\label{tab:teti_set1}
\begin{tabular}{lcc}
\hline\hline
Variable & Coefficient & (SE) \\
\hline
twitter\_sent (M1) & $-$0.012693 & (0.049061) \\
news\_sent (M3) & $-$0.122355*** & (0.031498) \\
twitter\_sent (M4) & $+$0.031176 & (0.049217) \\
news\_sent (M4) & $-$0.127532*** & (0.031600) \\
\hline\hline
\multicolumn{3}{p{0.9\linewidth}}{\footnotesize \textit{Notes:} HC3-robust SEs. Dependent variable: pct open-to-close return. *** $p<0.01$, ** $p<0.05$, * $p<0.1$.}
\end{tabular}
\end{table}

\textbf{SET 2 --- Panel FE with Binary Treatment.} Table~\ref{tab:teti_set2} reports the key heterogeneity coefficient $\beta_3$ from the interaction of lagged Twitter sentiment with a binary group indicator (\texttt{group = 1} if a firm's median market cap exceeds the cross-sectional median). None of the interaction terms are statistically significant at any lag. The signs alternate and the magnitudes are small relative to their standard errors. This is a null result: large-cap firms do not exhibit a measurably different sentiment--return relationship than small-cap firms in this sample.

\begin{table}[htbp]
\centering
\caption{Teti et al.\ (2019) Replication --- SET 2: Panel FE Binary Treatment ($\beta_3$ only)}
\label{tab:teti_set2}
\begin{tabular}{lcc}
\hline\hline
Variable & Coefficient & (SE) \\
\hline
ts\_lag\_0 $\times$ group & $+$0.058211 & (0.102712) \\
ts\_lag\_1 $\times$ group & $+$0.023445 & (0.089829) \\
ts\_lag\_2 $\times$ group & $-$0.051080 & (0.088592) \\
ts\_lag\_3 $\times$ group & $-$0.030345 & (0.086143) \\
ts\_lag\_4 $\times$ group & $+$0.029466 & (0.084418) \\
\hline\hline
\multicolumn{3}{p{0.9\linewidth}}{\footnotesize \textit{Notes:} Heteroskedasticity-robust SEs. Entity FE. \texttt{group = 1} if ticker median market cap $\geq$ cross-sectional median. \texttt{group} dropped from formula (collinear with entity FE); identification from interaction only. *** $p<0.01$, ** $p<0.05$, * $p<0.1$.}
\end{tabular}
\end{table}

\textbf{SET 3 --- Polarity-Weighted Index ($\text{ts\_b}$).} Table~\ref{tab:teti_set3} replaces the Bloomberg composite sentiment score with a custom polarity-weighted index $\text{ts\_b} = (\text{pos} - \text{neg}) / (\text{pos} + \text{neg} + \text{neu})$, which more precisely captures directional sentiment by normalizing for tweet volume. The lag-1 interaction term \texttt{tsb\_lag\_1} $\times$ \texttt{group} is marginally significant at $+$0.073 ($p < 0.10$). This is the first result in the replication to show any differentiation by firm size: larger firms appear to experience a slightly stronger positive return response to prior-day directional Twitter sentiment. However, the result is weak, does not persist beyond lag 1, and should be interpreted with caution.

\begin{table}[htbp]
\centering
\caption{Teti et al.\ (2019) Replication --- SET 3: Polarity-Weighted Index ($\text{ts\_b}$)}
\label{tab:teti_set3}
\begin{tabular}{lcc}
\hline\hline
Variable & Coefficient & (SE) \\
\hline
tsb\_lag\_0 $\times$ group & $+$0.060172 & (0.045303) \\
tsb\_lag\_1 $\times$ group & $+$0.072803* & (0.044166) \\
tsb\_lag\_2 $\times$ group & $-$0.044018 & (0.046582) \\
tsb\_lag\_3 $\times$ group & $-$0.030662 & (0.045037) \\
tsb\_lag\_4 $\times$ group & $-$0.014102 & (0.044444) \\
\hline\hline
\multicolumn{3}{p{0.9\linewidth}}{\footnotesize \textit{Notes:} $\text{ts\_b} = (\text{pos} - \text{neg}) / (\text{pos} + \text{neg} + \text{neu})$. Heteroskedasticity-robust SEs. Entity FE. *** $p<0.01$, ** $p<0.05$, * $p<0.1$.}
\end{tabular}
\end{table}

\newpage
\section{Gu \& Kurov (2020) Replication}

Gu \& Kurov (2020) use Fama-MacBeth (1973) cross-sectional regressions to examine whether lagged Twitter sentiment predicts open-to-open returns --- the appropriate return window given that Bloomberg releases Twitter sentiment at 9:20am, allowing execution at the 9:30am open. Their key finding is a significant positive predictive effect ($b \approx 0.136$) that passes a no-reversal test: only lag 1 is significant, consistent with information being fully absorbed within one trading day. I replicate their three core tests using risk-adjusted returns (Fama-French four-factor residuals), as the raw return specification encounters a rank-deficiency issue in the design matrix that is under investigation.

\textbf{Table 2 --- Return Predictability.} Table~\ref{tab:gk_t2} estimates the core Fama-MacBeth specification: open-to-open return on one-day lagged Twitter sentiment, controlling for five lags each of return, abnormal volume, Rogers-Satchell volatility, log market cap, and bid-ask spread. The risk-adjusted coefficient is $-$0.036 and is not statistically significant. This is directionally opposite to Gu \& Kurov's finding ($+$0.136***), and the magnitude is far smaller. The result suggests that in the current sample, lagged Twitter sentiment does not predict the next trading day's open-to-open return after controlling for known risk factors.

\begin{table}[htbp]
\centering
\caption{Gu \& Kurov (2020) Replication --- Table 2: Return Predictability}
\label{tab:gk_t2}
\begin{tabular}{lcc}
\hline\hline
 & (1) Raw Return & (2) Risk-Adj \\
\hline
twitter\_sent\_lag1 & --- & $-$0.035763 \\
 & --- & (0.049027) \\
\hline
Estimator & Fama-MacBeth & Fama-MacBeth \\
SE & Newey-West (5) & Newey-West (5) \\
\hline\hline
\multicolumn{3}{p{0.9\linewidth}}{\footnotesize \textit{Notes:} Fama-MacBeth. Newey-West SEs (5 lags). Open-to-open return. Controls: 5 lags each of return, abnorm\_vol, vol\_rs, log\_mkt\_cap, bid-ask spread. Raw return column omitted due to rank deficiency under investigation. *** $p<0.01$, ** $p<0.05$, * $p<0.1$.}
\end{tabular}
\end{table}

\textbf{Table 3 --- No-Reversal Test.} Table~\ref{tab:gk_t3} enters all five sentiment lags jointly. Under the Gu \& Kurov hypothesis, only lag 1 should be significant; significance at lags 3--5 would indicate reverse causality or momentum. All five lag coefficients are small and statistically insignificant. The inconclusive pattern here mirrors the no-reversal result from Section~4 (FE robustness): the sentiment signal, if it exists, cannot be isolated once lags are jointly conditioned on each other.

\begin{table}[htbp]
\centering
\caption{Gu \& Kurov (2020) Replication --- Table 3: No-Reversal Test}
\label{tab:gk_t3}
\begin{tabular}{lcc}
\hline\hline
 & (1) Raw Return & (2) Risk-Adj \\
\hline
twitter\_sent\_lag1 & --- & $-$0.019642 \\
 & --- & (0.051289) \\
twitter\_sent\_lag2 & --- & $+$0.007938 \\
 & --- & (0.058235) \\
twitter\_sent\_lag3 & --- & $-$0.042867 \\
 & --- & (0.061926) \\
twitter\_sent\_lag4 & --- & $-$0.055565 \\
 & --- & (0.059190) \\
twitter\_sent\_lag5 & --- & $-$0.017594 \\
 & --- & (0.054279) \\
\hline
Estimator & Fama-MacBeth & Fama-MacBeth \\
SE & Newey-West (5) & Newey-West (5) \\
\hline\hline
\multicolumn{3}{p{0.9\linewidth}}{\footnotesize \textit{Notes:} All lags entered jointly. Pass criterion: lag 1 significant, lags 2--5 insignificant (Gu \& Kurov, 2020). Result: inconclusive --- no lag is significant. *** $p<0.01$, ** $p<0.05$, * $p<0.1$.}
\end{tabular}
\end{table}

\textbf{Table 5 --- Twitter vs.\ News Sentiment.} Table~\ref{tab:gk_t5} compares the predictive power of Twitter and news sentiment, both separately and jointly. Neither channel is statistically significant in any specification. The joint model attenuates both coefficients slightly, consistent with a modest positive correlation between the two sentiment series. Gu \& Kurov report a contemporaneous correlation of approximately 0.20 between Twitter and news sentiment; the pattern of attenuation in the joint model is consistent with this, even though neither estimate clears conventional significance thresholds here.

\begin{table}[htbp]
\centering
\caption{Gu \& Kurov (2020) Replication --- Table 5: Twitter vs.\ News Sentiment}
\label{tab:gk_t5}
\begin{tabular}{lcc}
\hline\hline
 & (1) News Only & (2) Twitter + News \\
\hline
twitter\_sent\_lag1 & --- & $-$0.031707 \\
 & --- & (0.048244) \\
news\_sent\_lag1 & $-$0.023583 & $-$0.015916 \\
 & (0.047295) & (0.046690) \\
\hline
Estimator & Fama-MacBeth & Fama-MacBeth \\
SE & Newey-West (5) & Newey-West (5) \\
\hline\hline
\multicolumn{3}{p{0.9\linewidth}}{\footnotesize \textit{Notes:} Both sentiments use 1-day lag. Risk-adjusted open-to-open return. Low correlation implies largely orthogonal information channels. *** $p<0.01$, ** $p<0.05$, * $p<0.1$.}
\end{tabular}
\end{table}

\textbf{Table 6 --- Long-Short Trading Strategy.} Table~\ref{tab:gk_t6} reports the performance of a daily long-short strategy that goes long the top-decile and short the bottom-decile of Twitter sentiment each day, holding for 24 hours. The strategy generates a mean daily return of 0.066\%, translating to an annualized return of 16.70\% and a Sharpe ratio of 1.66. Gu \& Kurov report an annualized Sharpe of 3.17 over their 2015--2017 sample. The gap is notable: our Sharpe is roughly half theirs, though the strategy still produces economically meaningful returns before transaction costs. The win rate of 52.9\% indicates only marginal directional accuracy. One interpretation is that the sentiment signal retains some cross-sectional sorting power even when its time-series predictive coefficient is insignificant --- a pattern consistent with the Fama-MacBeth framework, where cross-sectional spread drives strategy returns independently of the average time-series slope.

\begin{table}[htbp]
\centering
\caption{Gu \& Kurov (2020) Replication --- Table 6: Long-Short Trading Strategy}
\label{tab:gk_t6}
\begin{tabular}{lc}
\hline\hline
Statistic & Value \\
\hline
Mean daily L-S return (\%) & 0.0663 \\
Annualized return (\%) & 16.70 \\
Annualized Sharpe & 1.66 \\
Win rate & 52.9\% \\
Trading days & 329 \\
\hline\hline
\multicolumn{2}{p{0.7\linewidth}}{\footnotesize \textit{Notes:} Long (short) = top (bottom) decile daily Twitter sentiment. 24-hour hold. Before transaction costs. Gu \& Kurov (2020) report Sharpe of 3.17 over Jan 2015--Feb 2017 (537 days).}
\end{tabular}
\end{table}

\newpage
\section{Summary and Next Steps}

The results across all specifications and replications tell a consistent story: Twitter sentiment, as measured by Bloomberg's composite daily score, does not reliably predict stock returns in this sample. The FE and DoubleML estimates in Sections 2--4 find no significant effect. The Teti et al.\ replication finds null results in SET 1 and SET 2, with only a marginal signal in SET 3 (the polarity-weighted $\text{ts\_b}$ index interacted with firm size at lag 1, $p < 0.10$). The Gu \& Kurov replication finds no significant return predictability in the Fama-MacBeth framework, and a long-short Sharpe ratio of 1.66 --- economically interesting but well below the 3.17 reported in the original paper.

News sentiment, by contrast, consistently emerges as a significant predictor: it is significant in FE Model 2, in the Teti pooled OLS, and in the DoubleML PLR. This asymmetry between Twitter and news is meaningful. It may reflect a decline in Twitter's role as an information source since the 2015--2017 period studied by Gu \& Kurov, or it may reflect the noisier nature of social media relative to structured news feeds. It may also be that Bloomberg's composite Twitter sentiment score, which aggregates across all tweets for a firm in a day, obscures the signal that would be present in more granular or entity-specific tweet data.

Several avenues remain open. First, the raw return rank-deficiency issue in the Gu \& Kurov Fama-MacBeth models warrants a targeted diagnostic. Second, DoubleML robustness checks for both replications (lagged treatments, impact persistence, no-reversal joint tests) are underway. Third, and most substantively, a custom sentiment analysis tool is in development that will enable filtering by specific individuals or entities and precise tweet timing --- an approach that may recover the causal signal that composite daily scores appear to obscure.

\end{document}
